\section{The determinant energy and its boundary flux}
\label{sec:energy}

The time component of the wave current is a sum of squares. After
regularization it gives the density for a spatial Dirichlet potential:
a lower bound on the density throughout a ball forces a definite
depth at its center. To compare that depth at different times, we
need a scalar energy on the three real coordinates
$z=(T,\tau)\in\RR\times\CC_{\mathbf k}$ whose Laplacian is bounded
above by the wave flux through the side of the spatial fiber.

Keep the annular wave on the torus from Section~\ref{sec:current},
normalized by $\|f\|_2=1$. Put $w_0=1/10$ and assume its sector
$|v|\le\gamma$ has $\gamma<w_0/4$.
Write $s=|\tau|$ and $a=Kh$, where $K\ge1$ is fixed and $h$ is small
enough that $4a\le1$. All the constructions of this section take place
on $a<s<4$.

\subsection{The geometry of the fibers}

At fixed $z$, let $X$ run over the spatial torus and let the imaginary
spatial variables $(P,Z)\in\CC_{\mathbf k}^d\times\CC_{\mathbf k}$ vary
in the domain
\begin{equation}\label{eq:fiber-domain}
 \Omega_z:\quad r=\frac{|W|^2}{w_0^2}+100|Q|^2-1<0,
 \qquad Q=\Phi/a.
\end{equation}
The shape of the fiber reflects two properties of the current from
Proposition~\ref{prop:current}: it is tangent to the level sets of
$\Phi$, and its Gaussian is centered at $W=v$ for a mode of slope $v$.
Only the $W$ part of the wall normal can therefore carry wave flux.
In $(W,Q)$ the fibers are all the same ellipsoid. The definitions
$Y=(\tau+Z)/\sqrt2$, $W=P/Y$ and
$\Phi=(Z-\tau)/\sqrt2+P\cdot P/(2Y)$ give the inverse map
\begin{equation}\label{eq:fiber-inverse}
 Y=\frac{\sqrt2\,\tau+aQ}{1+W\cdot W/2},\qquad P=WY,\qquad
 Z=\sqrt2\,Y-\tau .
\end{equation}
The denominator is bounded away from zero on $|W|\le w_0$, and
$|aQ|\le s/10$, so $y=|Y|\asymp s$. Both the chart and its inverse
are holomorphic on neighborhoods of these closed fibers and depend
smoothly on $\tau$. On the total domain
$\Omega=\{(z,X,P,Z):a<s<4,\ (X,P,Z)\in\Omega_z\}$ we use $H,G$ for
the full Hessian and gradient from Lemma~\ref{lem:hessian-calculus}.
Their spatial parts at fixed $z$ are $H_{\rm fib},G_{\rm fib}$. Fiber
integrals use ordinary volume $\dd\zeta=\dd X\,\dd P\,\dd Z$ in the
original coordinates.

\begin{lemma}[Fiber geometry]\label{lem:fiber-geometry}
The spatial Hessian $H_{\rm fib}(r)$ is positive definite. The full
$H(r)$ is positive semidefinite with a one-dimensional quaternionic
kernel contained in both $G(r)^\perp$ and $\ker H(|\Phi|^2)$. On the
closed fibers, uniformly for $a<s<4$,
\begin{align}
 H_{\rm fib}(r)&\ge c s^{-2}I,\qquad
 G_{\rm fib}(r)G_{\rm fib}(r)^*\le H_{\rm fib}(r),\label{eq:fiber-curvature}\\
 \dd P\,\dd Z&\asymp a^2s^{2d}\dd W\,\dd Q,\qquad
 \det H_{\rm fib}(r)\asymp a^{-2}s^{-2d}.\label{eq:fiber-jacobian}
\end{align}
The spatial $H_{\rm fib}(|\Phi|^2)$ has rank one with nonzero eigenvalue
comparable to one. Also $|G(|W|^2)|\lesssim s^{-1}$.
\end{lemma}

\begin{proof}
Set $F=(W/w_0,10Q)$ and let $J$ be its complex Jacobian in
$(\tau,P,Z)$, with spatial submatrix $J_{\rm fib}$. Since $r=|F|^2-1$,
Lemma~\ref{lem:hessian-calculus} gives
\[
 H(r)=4J^*J,\qquad G(r)=2J^*F\mathbf i .
\]
At fixed $\tau$ the spatial Jacobian $J_{\rm fib}$ is invertible.
In the full space we can vary $\tau$ while keeping $(W,Q)$ fixed;
the resulting direction spans $\ker J$ over $\CC_{\mathbf k}$ and
its quaternionic span is the null line of $H(r)$. This line is
orthogonal to $G(r)$, which belongs to the image of $J^*$, and lies in
$\ker H(|\Phi|^2)$ because the last row of $J$ is proportional to
$\partial\Phi$. By
\eqref{eq:fiber-inverse}, derivatives of $(P,Z)$ in $W$ are $O(s)$
and those in $Q$ are $O(a)$. Thus $\|J_{\rm fib}^{-1}\|\lesssim s$;
together with $|F|\le1$ this proves \eqref{eq:fiber-curvature}.

The complex Jacobian determinant of \eqref{eq:fiber-inverse} in $(W,Q)$
is $\sqrt2\,aY^d/(1+W\cdot W/2)$: taking the exterior product of the
differentials of $P=WY$ with $\dd Y$ leaves just $Y^d$ times the product
of the differentials of $W$ with $\dd Y$. Squaring the absolute value
gives the volume formula. The determinant of $4J_{\rm fib}^*J_{\rm fib}$
is $4^n|\det_{\CC}J_{\rm fib}|^2$, giving the second bound in
\eqref{eq:fiber-jacobian}. Finally the spatial derivatives of $\Phi$ are
bounded and $\partial_Z\Phi=(1-W\cdot W/2)/\sqrt2$ is bounded away from
zero. All first derivatives of $W=P/Y$ are $O(s^{-1})$.
\end{proof}

\subsection{The Dirichlet potential}

The time density can vanish. We first smooth the wave current in $X$
and cut it off in the passive coordinate $Q$. To make its time
density strictly positive, we then add a conserved current whose image
is the null line of $H(r)$; it has zero flux through the wall.
For $\ell>0$ let $J_\ell$ be the periodization on the spatial torus of
the probability kernel
$c_n\ell^{-n}(1+|X/\ell|^2)^{-n-2}$. Choose a smooth $0\le g\le1$ on
$\CC_{\mathbf k}$, equal to one on $|Q|\le1/100$ and zero on
$|Q|\ge1/20$. Define
\begin{equation}\label{eq:regularized-current}
 A_{\rm w}=g(Q)(J_h*_X A^0),\qquad
 A=A_{\rm w}+h^{50n}C(H(r)),\qquad \mu=A_{TT}.
\end{equation}
Convolution takes place only in the real spatial variable. Since
$\partial\Phi$ is independent of $X$, the averaged current still
annihilates the image of $H(|\Phi|^2)$ and hence $G(g(Q))$. The product
rule for divergence therefore preserves conservation under the cutoff.
By Lemmas~\ref{lem:quaternionic-matrices} and~\ref{lem:cofactor-divergence},
the regularizing cofactor is positive semidefinite and divergence free.
Its image is $\ker H(r)$, perpendicular to $G(r)$, and its $TT$ entry
is $\det H_{\rm fib}(r)>0$. Thus
\begin{equation}\label{eq:current-background-orthogonality}
 A\ge0,\qquad \mu>0,\qquad A:H(|\Phi|^2)=0 .
\end{equation}
Put $B_0=\lambda|\Phi|^2$, $B=H(B_0)$, $B_{\rm fib}=H_{\rm fib}(B_0)$,
with $\lambda=h^{-30n}$. By \eqref{eq:current-background-orthogonality},
the full background satisfies $A:B=0$. Its spatial block $B_{\rm fib}$
has one nonzero eigenvalue $\beta\asymp\lambda$. The factor
$h^{50n}$ keeps the added density below the wave estimates
of Section~\ref{sec:barrier}; we take $\lambda$ large enough that
\begin{equation}\label{eq:energy-density-ceiling}
 \mu^{1/d}\le\lambda
\end{equation}
for all sufficiently small $h$. To see this, each Gaussian amplitude
factor in $\chi$ has size at most $(\kappa/y)^{d/2}\lesssim h^{-d}$.
The columns $V^v$ are bounded and $\sum|c_{\kappa,v}|\lesssim h^{-n/2}$
by the cap estimate in Section~\ref{sec:current}, so
$\|A^0\|\lesssim h^{-n-2d}$. The regularizer adds only
$O(h^{50n}a^{-2}s^{-2d})$ to $\mu$ by \eqref{eq:fiber-jacobian}.

On each spatial fiber there is a unique smooth real $u$ satisfying
\begin{equation}\label{eq:fiber-equation}
 D_u=H_{\rm fib}(B_0+u)>0,\qquad
 \det D_u=\lambda\mu,\qquad u|_{\partial\Omega_z}=0 .
\end{equation}
Indeed, the spatial defining Hessian is positive definite, and
$B_{\rm fib}\ge0$ has zero determinant since $n\ge2$. Together with
$\mu>0$ these are the hypotheses of Theorem~\ref{thm:dirichlet}.
The chart \eqref{eq:fiber-inverse} and the parameter assertion of that
theorem give smooth dependence on $z$ up to the lateral boundary;
its barriers \eqref{eq:dirichlet-barriers} give $u\le0$.
To see how the equation measures depth, add $\alpha I$ to the spatial
background. Then
$\det(B_{\rm fib}+\alpha I)=(\beta+\alpha)\alpha^d$; thus for
$0<\alpha\le\lambda$ a density $\lambda\mu$ corresponds to an
isotropic comparison quadratic of curvature $\alpha\asymp\mu^{1/d}$.
Section~\ref{sec:barrier} will use such quadratics on balls.

\subsection{Cancelling the full determinant}

The quantity that measures the depth of the potential against its
density is $X_0(z)=\int_{\Omega_z}(-u)\mu\dd\zeta$. Since the real
columns of $A$ have zero divergence, Green's formula for $\Delta_z X_0$
has interior term $-\int_{\Omega_z}A:H(u)\dd\zeta$. In a related
estimate for real positive tensors, Serre uses the determinant of a
Schur complement to isolate a diagonal entry of a second tensor
\cite[Proposition~2.2]{Serre2023Mixed}. Here only the spatial block of
$H(B_0+u)$ is known to be positive. Separate off its scalar Schur
complement:
\begin{equation}\label{eq:energy-schur}
 H(B_0+u)=
   \begin{pmatrix}\ell&c^*\\ c&D_u\end{pmatrix}
   =\sigma_0 e_Te_T^*+\Sigma,\qquad
 \sigma_0=\ell-c^*D_u^{-1}c,\qquad
 \Sigma=\begin{pmatrix}c^*D_u^{-1}c&c^*\\ c&D_u\end{pmatrix}\ge0 .
\end{equation}
Here $e_T$ is the time coordinate vector. Completing the square gives
$\Sigma\ge0$. In the contraction with $A$ the background term $A:B$
vanishes by \eqref{eq:current-background-orthogonality}; in the
determinant, $\det D_u=\lambda\mu$ gives
$\det H(B_0+u)=\lambda\mu\sigma_0$. Thus
\begin{equation}\label{eq:schur-energy-identity}
 A:H(u)=A:\Sigma+\mu\sigma_0
        =A:\Sigma+\lambda^{-1}\det H(B_0+u).
\end{equation}
The full determinant has no prescribed sign. We look for a second
fiber integral whose Laplacian cancels this term. Write
$D_v=H_{\rm fib}(B_0+v)$. Along the segment $B_0+\theta u$, the
cofactor of the full Hessian has $TT$ entry $\det D_{\theta u}$, and
its contraction with $H(u)$ differentiates the full determinant in
$\theta$. This suggests integrating the spatial determinant along the
segment. For a smooth real $v$ on $\Omega_z$ with $D_v>0$ and zero
boundary values, set
\begin{equation}\label{eq:energy-functional}
 E_z(v)=\frac1\lambda\int_{\Omega_z}v
            \int_0^1\det D_{\theta v}\dd\theta\,\dd\zeta,\qquad
 \mathcal J_z(v)=E_z(v)-\int_{\Omega_z}v\mu\dd\zeta .
\end{equation}
The segment integral is the quaternionic counterpart of the
Aubin--Mabuchi energy: Mabuchi's functional $L$ integrates the variation
of a Kähler potential against its complex Monge--Ampère measure
\cite[(2.2.1)]{Mabuchi1986}. Here the determinant is a polynomial,
so the integrand is smooth up to the singular endpoint $D_0=B_{\rm fib}$.
Define $\mathcal F(z)=\mathcal J_z(u)=E_z(u)+X_0(z)$.

\begin{proposition}[Energy balance]\label{prop:energy-flux}
The energy $\mathcal F$ is smooth on $a<|\tau|<4$. Let $\rho$ denote
the boundary value of $u/r$, which extends smoothly from the interior.
In the three variables $z=(T,\tau)$ one has
\begin{equation}\label{eq:energy-flux}
 \Delta_z\mathcal F(z)+\int_{\Omega_z} A:\Sigma\,\dd\zeta
   =\int \rho\,G(r)^*A_{\rm w}G(r)\,\delta(r)\dd\zeta .
\end{equation}
The normal seen by the current is that of the moving wall in the
total space, so $G(r)$ is the \emph{full} gradient. The integration is
on one fiber: $\delta(r)\dd\zeta=\dd S_{\rm fib}/|\nabla_{\rm fib}r|$,
where the spatial gradient is nonzero by the ellipsoid chart. Since
$A:\Sigma\ge0$, the boundary integral bounds $\Delta_z\mathcal F$
from above.
\end{proposition}

\begin{proof}
Denote the boundary integral in \eqref{eq:energy-flux} by
$\mathcal B(z)$. The two identities whose sum gives the balance are
\begin{align*}
  \Delta_zX_0
      &=\mathcal B(z)-\int_{\Omega_z}A:\Sigma\dd\zeta
                       -\lambda^{-1}\int_{\Omega_z}\det H(B_0+u)\dd\zeta,\\
  \Delta_z[E_z(u)]
      &=\lambda^{-1}\int_{\Omega_z}\det H(B_0+u)\dd\zeta .
\end{align*}
For a fixed complex product domain, \citet[Theorem~3.1]{AhagEtAl2009}
identified the Laplacian of the slice energy, up to normalization,
with the integral of the full Monge--Ampère measure over the fiber.
Here the fiber moves, so the second identity requires a boundary calculation.
Both follow by testing with a smooth $\eta=\eta(z)$ compactly supported in the
parameter domain. The only entry of $H(\eta)$ is
$H_{TT}(\eta)=\Delta_z\eta$. For any Hermitian matrix field $M$ with
divergence-free real columns, Green's formula \eqref{eq:calculus-green}
on the \emph{total} domain gives
\begin{equation}\label{eq:green-moving-fiber}
 \int_\Omega u M_{TT}\Delta_z\eta
 =\int_\Omega \eta\,M:H(u)
   -\int_{\partial_{\rm lat}\Omega}
       \eta\,(\nu^*M\nu)\,\partial_{\mathbf n}u\dd S .
\end{equation}
Here $\partial_{\rm lat}\Omega$ is the side $r=0$, and volume on
$\Omega$ is $\dd z\,\dd\zeta$. The term containing
$u\nabla\eta$ vanishes because $u=0$ on the lateral boundary; there
are no end terms because of the support of $\eta$. The unit outward
normal $\mathbf n=\nabla r/|\nabla r|$ on the total domain has quaternionic
column $\nu=G(r)/|\nabla r|$. Since $r$ is a smooth defining function,
$u=r\rho$ near the boundary with $\rho$ smooth and $\rho\ge0$ on $r=0$.
In particular $\partial_{\mathbf n}u=\rho|\nabla r|$.

For $X_0$, take $M=A$ in \eqref{eq:green-moving-fiber} and change
sign since $X_0$ integrates $(-u)\mu$. The interior term follows
from \eqref{eq:schur-energy-identity}. On the wall the regularizer
$C(H(r))$ has image perpendicular to $\nu$, so only $A_{\rm w}$
contributes. Substituting $\partial_{\mathbf n}u=\rho|\nabla r|$ and
$\dd S=|\nabla r|\delta(r)\dd z\,\dd\zeta$ gives $\mathcal B(z)$.

For $E_z(u)$, take $M=C(H(B_0+\theta u))$, whose columns are conserved
by Lemma~\ref{lem:cofactor-divergence}. Here the entire normal
coefficient vanishes. Indeed \eqref{eq:boundary-restriction} gives,
on $r=0$,
\[
 H(B_0+\theta u)|_{\nu^\perp}
       =(B+\theta\rho H(r))|_{\nu^\perp}.
\]
The null line of $H(r)$ belongs to both $\nu^\perp$ and $\ker B$ by
Lemma~\ref{lem:fiber-geometry}; the restricted form is therefore
singular. The cofactor restriction formula in
Lemma~\ref{lem:quaternionic-matrices} yields $\nu^*M\nu=0$, independently
of the sign of the full Hessian. The same formula in the $e_T$ direction
gives $M_{TT}=\det D_{\theta u}$. Now integrate
\eqref{eq:green-moving-fiber} in $\theta$ and divide by $\lambda$.
On the left we obtain $\int E_z(u)\Delta_z\eta\,\dd z$; on the right,
integration of $M:H(u)=\partial_\theta\det H(B_0+\theta u)$ gives the
full determinant at $\theta=1$, since $\det B=0$. This proves the
second identity and hence the balance in distributions. The chart
\eqref{eq:fiber-inverse} makes the volume and surface integrals smooth
in $z$, so it also holds pointwise.
\end{proof}

\subsection{Size of the energy on a fiber}

Fix $z$. The segment integral in \eqref{eq:energy-functional} is a
primitive of the spatial determinant: for a smooth variation $w$ with
zero boundary values, one has
\begin{align*}
 \mathrm d\mathcal J_z(v)[w]
   &=\int_{\Omega_z}w(\lambda^{-1}\det D_v-\mu)\dd\zeta,\\
 \mathrm d^2\mathcal J_z(v)[w,w]
   &=-\lambda^{-1}\int_{\Omega_z}
       G_{\rm fib}(w)^*C(D_v)G_{\rm fib}(w)\dd\zeta .
\end{align*}
For the first identity, differentiating under the integral in
\eqref{eq:energy-functional} produces a term
$\theta v C(D_{\theta v}):H_{\rm fib}(w)$. By cofactor divergence one
may interchange $v$ and $w$ in its fiber integral; both vanish on the
boundary. The integrand in $\theta$ is then
$w\,\partial_\theta(\theta\det D_{\theta v})$. Differentiating once
more and integrating by parts gives the second formula. It is strictly
negative for every nonzero $w$, since $C(D_v)>0$. Thus the Dirichlet
solution is the unique maximizer on this convex class. This is the
smooth determinant variational principle with a rank-one background.
For zero background, Wan proved the first variation and the corresponding
variational characterization in quaternionic finite energy classes
\cite[Section~4]{Wan2020}.

Since $u\le0$ and $D_{\theta u}=(1-\theta)B_{\rm fib}+\theta D_u\ge0$,
the correction $E_z(u)$ is nonpositive. At a solution the linear
term in $\mathcal F$ uses the determinant at $B_0+u$. Thus
$\lambda\mathcal F$ has the form of Aubin's $J$-functional
\cite[p.~145]{Aubin1984}, with $B_0+u$ as reference potential and $B_0$
as argument in the relative formula \cite[Section~2.1]{BBGZ2013}.
The rank-one background lets us calculate how much of $X_0$ remains
in $\mathcal F=X_0+E_z(u)$.

\begin{lemma}[Energy comparison]\label{lem:energy-comparison}
On every spatial fiber,
\begin{equation}\label{eq:energy-exact}
 \mathcal F=\frac{n-1}{n+1}X_0+
   \frac{1}{\lambda n(n+1)}
   \int_{\Omega_z}G_{\rm fib}(u)^*C(D_u)G_{\rm fib}(u)\dd\zeta .
\end{equation}
In particular,
$\frac{n-1}{n+1}X_0\le\mathcal F\le\frac{n}{n+1}X_0$.
\end{lemma}

\begin{proof}
Because $B_{\rm fib}$ has rank one, the determinant along the segment
has just two terms:
\[
 \det D_{\theta u}=\theta^n\det D_u+
    (1-\theta)\theta^{n-1} C(D_u):B_{\rm fib}.
\]
Write $I=\int_{\Omega_z}(-u) C(D_u):B_{\rm fib}\dd\zeta\ge0$.
Integrating in $\theta$ gives
$\mathcal F=nX_0/(n+1)-I/[\lambda n(n+1)]$, in particular the upper
bound. Also $C(D_u):D_u=n\det D_u$, so
\[
 n\lambda X_0-I
   =\int_{\Omega_z}(-u)C(D_u):H_{\rm fib}(u)\dd\zeta
   =\int_{\Omega_z}G_{\rm fib}(u)^*C(D_u)G_{\rm fib}(u)\dd\zeta .
\]
This last integration by parts uses $u=0$ at the boundary. Substitution
proves \eqref{eq:energy-exact} and its lower bound.
\end{proof}
